\scanpage{277}
\mspage{18}

Ce foncteur s'identifie au foncteur de composition
\[ c\longmapsto c\circ\alpha \]
où $\alpha$ est le foncteur d'inclusion
\[ \left\{\begin{aligned}
\varepsilon&\xrightarrow{\alpha} I\\
\varepsilon&=\{0,1\},\quad\text{catégorie discrète à deux objets}\\
&\alpha\ \text{donné par}\quad \alpha(0)=\text{source de }I\\
&\phantom{\alpha\ \text{donné par}\quad}\alpha(1)=\text{but de }I
\end{aligned}\right. \]

\unsure{$\rightarrow$|,} Si $\omega'$ est l'\deleted{\unsure{caté}}orientation opposée de $I$, on a un diagramme commutatif
\diagram{\begin{tikzcd}[ampersand replacement=\&,row sep=large,column sep=large]
\underline{\mathrm{Ch}}_{(I,\omega)}(X) \arrow[r,dash,"{\varphi}"{yshift=4pt},"\sim"] \arrow[d,"{(\unsure{\sigma|\alpha},\beta)}"] \& \underline{\mathrm{Ch}}_{(I,\omega')}(X) \arrow[d,dash,"{(\sigma',\beta')}"]\\
X\times X \arrow[r,"s"] \& X\times X
\end{tikzcd}}
% slanted
\margin{Le foncteur $\varphi$ est celui qui associe à tout chemin $c$, de type $\tau$, le \uline{chemin opposé} $c^{\mathrm{opp}}$ ou $c^{-1}$, de type $\tau^{\mathrm{opp}}$. Si $\tau=\tau^{\mathrm{opp}}$, c'est donc un chemin \uline{\unsure{du|de} même type} $\tau$, \deleted{\unsure{l'aut}} $\varphi$ est un \uline{\unsure{autom}} de $\underline{\mathrm{Ch}}_{\tau}(X)_{\unsure{s}}$. Cet \unsure{autom} \inserted{(\unsure{naturel en|variant en} $X$)} est déduit d'un \unsure{autom.} $\theta$ de $I$, cf ci-dessous.}
$s$ est le foncteur de symétrie
\[ s(x,y)=(y,x). \]

Donc l'isomorphisme $\varphi$ ne respecte pas les flèches structurales ``source-but'' $(\sigma,\beta)$.

Il y a cependant un \unsure{cas heureux|mieux} quand il existe un iso. d'intervalles \uline{orientés}
\[ \theta\colon(I,\omega)\xrightarrow{\sim}(I,\omega'), \]
i.e. quand $I$ est \uline{symétrique}. Ce n'est pas le cas pour les $\mathbb{\Delta}_i$ $(i\geq 1)$ (ces catégories n'ont pas d'autre \unsure{autom.} que l'identité), mais c'est le cas pour les catégories $\Phi_1,\Psi_1$ de l'exemple $4^{\mathrm{o}}$, plus généralement \deleted{pour les}
% archive sheet number: 276
